\DocumentMetadata{
  pdfversion=2.0,
  pdfstandard=ua-2,
  lang=en-US,
  tagging=on,
  tagging-setup={math/setup=mathml-SE,tabsorder=structure}
}
\documentclass[11pt,letterpaper]{article}
\usepackage[margin=0.68in,includefoot]{geometry}
\usepackage{newcomputermodern}
\usepackage{amsmath,amscd,mathtools}
\providecommand{\square}{\mdlgwhtsquare}
\usepackage[hidelinks]{hyperref}
\hypersetup{
  pdftitle={Analysis First-Year Exam, August 2023, Part 2},
  pdfauthor={University of Florida Department of Mathematics},
  pdfsubject={Graduate mathematics examination},
  pdfdisplaydoctitle=true,
  bookmarksopen=true
}
\setcounter{secnumdepth}{0}
\setlength{\parindent}{0pt}
\setlength{\parskip}{0.28em}
\setlength{\emergencystretch}{3em}
\setlength{\leftmargini}{2.15em}
\setlength{\leftmarginii}{2.65em}
\setlength{\leftmarginiii}{3em}
\renewcommand{\labelenumi}{\textbf{\theenumi.}}
\pagestyle{empty}
\raggedbottom
\allowdisplaybreaks
\newenvironment{examproblems}[1]{%
  \begin{enumerate}%
  \setcounter{enumi}{#1}%
  \setlength{\topsep}{0.35em}%
  \setlength{\partopsep}{0pt}%
  \setlength{\itemsep}{0.48em}%
  \setlength{\parsep}{0.18em}%
}{\end{enumerate}}
\newenvironment{examsubparts}{%
  \begin{enumerate}%
  \setlength{\topsep}{0.18em}%
  \setlength{\partopsep}{0pt}%
  \setlength{\itemsep}{0.16em}%
  \setlength{\parsep}{0.1em}%
}{\end{enumerate}}
\newenvironment{examinstructions}{%
  \par\begingroup\itshape
}{%
  \par\endgroup\vspace{0.18em}
}
\makeatletter
\renewcommand\section{\@startsection{section}{1}{0pt}%
  {-1.25ex plus -.25ex minus -.1ex}%
  {0.55ex plus .1ex}%
  {\normalfont\large\bfseries}}
\renewcommand{\@maketitle}{%
  \newpage\null\vskip -1.2em%
  {\footnotesize\bfseries
    \href{https://www.ufl.edu/}{University of Florida}, \href{https://math.ufl.edu/}{Department of Mathematics}\par}%
  \vskip 0.18em%
  \tagstructbegin{tag=Title}%
    {\LARGE\bfseries\href{https://gma.math.ufl.edu/exams/analysis-first-year/}{Analysis First-Year Exam}, August 2023, Part 2\par}%
  \tagstructend%
  \vskip 0.35em%
  \tagmcbegin{artifact}%
    \hrule height 1pt%
  \tagmcend%
  \vskip 0.55em%
}
\makeatother
\title{Analysis First-Year Exam, August 2023, Part 2}
\author{}
\date{}
\begin{document}
\maketitle
\thispagestyle{empty}
\begin{examinstructions}
Answer FOUR questions in detail. State carefully any results used without proof.
\end{examinstructions}
\begin{examproblems}{0}
\item
Let \((f_n)\) be a sequence of continuous functions from \([0,1]\) to \([0,1]\). For each of the following, give a proof (if true) or a counterexample (if false):
\begin{examsubparts}
\item[\textnormal{(i)}]
if \(f_n(t)\) decreases pointwise to~\(0\) then \(\int_0^1 f_n(t)\,dt \to 0\);
\item[\textnormal{(ii)}]
if \(f_n(t)\) only converges pointwise to~\(0\) then \(\int_0^1 f_n(t)\,dt \to 0\).
\end{examsubparts}
\item
For \(0<t<1\) and for~\(n\) a positive integer, let \(f_n(t)=t^n\).
\begin{examsubparts}
\item[\textnormal{(i)}]
Show that the sequence \((f_n)\) converges uniformly on each compact subset of \((0,1)\).
\item[\textnormal{(ii)}]
Show that \((f_n)\) does not converge uniformly on \((0,1)\).
\end{examsubparts}
\item
Let the function \(f:[1,\infty)\to\mathbb{R}\) be continuous, with \(\lim_{t\to\infty}f(t)=A\in\mathbb{R}\). Prove there exists a sequence \((p_n)\) of polynomials such that \(p_n(1/t)\to f(t)\) uniformly for \(t\geq 1\).
\item
Let \(A\subseteq[0,1]\) have positive Lebesgue measure: that is, \(\lambda(A)>0\). Prove that if \(0<\alpha<\lambda(A)\) then there exists a Lebesgue measurable set \(A_\alpha\subseteq A\) such that \(\lambda(A_\alpha)=\alpha\).
\item
State the Monotone Convergence Theorem and ‘Fatou’s Lemma’. Deduce one of these from the other.
\end{examproblems}
\end{document}
